% VI/46 reader bridge. Additive exposition; all original exercises retained.
\paragraph{Restriction surjectivity is stronger than gluing.}
A sheaf $\mathcal F$ of abelian groups on $X$ is \emph{flabby}
(or flasque) when for every pair of open sets $V\subseteq U$
the restriction $\mathcal F(U)\to\mathcal F(V)$ is surjective.
The ordinary sheaf axiom says that compatible local sections
\emph{glue}; flabbiness additionally says that every section on
$V$ \emph{extends} to all of $U$.  This is why flabbiness is
special and why its cohomological consequences are strong.

\paragraph{Worked example 1: the skyscraper sheaf.}
Choose $p\in X$ and an abelian group $G$.  For the skyscraper
sheaf $i_*G$ with $i:\{p\}\hookrightarrow X$, its sections are
\[
 (i_*G)(U)=
 \begin{cases}G,&p\in U,\\0,&p\notin U.\end{cases}
\]
Every restriction is either an identity on $G$, the map
$G\to0$, or $0\to0$, and is therefore surjective.  Thus
$i_*G$ is flabby.  Notice that a skyscraper sheaf need not
look like a constant sheaf on all of $X$: its section group
changes when an open set stops containing the distinguished point.

\paragraph{Worked example 2: all functions versus locally constant functions.}
For any abelian group $G$, let $\mathcal A(U)$ be the group of
\emph{all} set-theoretic functions $U\to G$.  A function on
$V\subseteq U$ extends by assigning $0\in G$ to every point
of $U\setminus V$, hence $\mathcal A$ is flabby.  This extension
need not be continuous; continuity is not part of this sheaf.
On the other hand, take $X=\mathbb R$ with the ordinary topology
and let $\underline{\mathbb Z}$ be the sheaf of locally constant
integer-valued functions.  On the disconnected open set
$V=(-2,-1)\cup(1,2)$, the section equal to $0$ on the first
component and $1$ on the second cannot extend to all of the
connected space $\mathbb R$ as a locally constant function.
Therefore $\underline{\mathbb Z}$ is not flabby.  (The topology
is specified here; on other spaces the answer can differ.)

\paragraph{Why a flabby sheaf has no higher cohomology.}
Flabbiness provides global lifting of local data: for a short exact
sequence of sheaves of abelian groups
\[
 0\longrightarrow\mathcal F\longrightarrow\mathcal G
 \longrightarrow\mathcal H\longrightarrow0,
\]
if $\mathcal F$ is flabby, then
$\Gamma(X,\mathcal G)\to\Gamma(X,\mathcal H)$ is surjective.
Indeed, locally lift a section of $\mathcal H$; on overlaps the
differences lie in $\mathcal F$, and the extension property
allows a maximal partial lifting to be enlarged until it is
global.  The maximal-extension argument is the missing step in
trying to deduce right exactness of global sections from mere
stalkwise exactness.  In particular, every flabby sheaf is
acyclic for the global-section functor:
\[
 H^q(X,\mathcal F)=0\qquad(q>0).
\]
For a sheaf of abelian groups this is ordinary sheaf cohomology;
it is not a statement that every sheaf is flabby or that
restriction maps in every exact sequence are surjective.

\paragraph{Worked example 3: flabbiness survives direct image.}
For a continuous map $f:X\to Y$, suppose $\mathcal F$ is flabby
on $X$.  If $V\subseteq U$ are open in $Y$, then
\[
 (f_*\mathcal F)(U)=\mathcal F(f^{-1}U)
 \longrightarrow\mathcal F(f^{-1}V)=(f_*\mathcal F)(V)
\]
is surjective because $f^{-1}V\subseteq f^{-1}U$.  Hence
$f_*\mathcal F$ is flabby.  This calculation is elementary,
but illustrates the right way to check permanence properties:
write the actual restriction map rather than appealing only
to the word ``pushforward.''

\paragraph{Bridge to \v{C}ech cohomology.}
Flabby sheaves also have vanishing positive-degree \v{C}ech
cohomology for every open cover.  The next section displays the
actual cochain differential, then contrasts a flabby example
with a nontrivial class of $\mathcal O_{\mathbb P^1}(-2)$.
These statements should not be conflated with the claim that
arbitrary \v{C}ech cohomology automatically equals derived-functor
sheaf cohomology without hypotheses on the cover and sheaf.
